\titledquestion{Run factor graph SLAM on simulated data}
\label{q:task02}
You are given a data set consisting of
\begin{itemize}
    \item odometry (K, 3): the measured movement of the robot in body frame
    \item z (K,): Python list of detections at each time step. Each detection is (\#measurements, 2).
    \item poseGT (K, 3): the pose ground truth at each time step
    \item landmarks (M, 2): the ground truth landmarks
\end{itemize}

We have given you a set of pretty good initial tuning values (marked with \pythoninline{# TODO tune} in \pythoninline{configs/sim_default.yaml}). We want you to find 3 different sets of tuning values that differ in different ways and explain the results you get using them.

An example of three different sets of tuning values could be
\begin{itemize}
    \item a set that gives better performance than the initial tuning values in some way
    \item a set that gives similar performance in terms of error but bad performance in terms of consistency
    \item a set that yields very few landmarks, but still manages to keep track to some degree
\end{itemize}

\hfill

In addition to the plots listed on the first page, we want one figure that has no counterpart in the EKF-SLAM assignment: for a pose \(x_k\) somewhere in the middle of the run, compare its error \emph{as estimated at time \(k\)} with its error \emph{in the final estimate}.
A filter cannot revise the past; a smoother can, and does. Explain what you see, and relate it to (9.2) and to what the odometry factors are doing to the earlier poses when a landmark is re-observed.

\hfill

Some questions to help discussion (not mandatory to answer):
\begin{itemize}
    \item Which types of errors do you see?
    \item Why do you think the errors are like this?
    \item Setting \pythoninline{association: method: gt} runs with perfect data association. How much of your error is the front-end responsible for, and how much is the back-end? (Sec.\ 9.1)
    \item Would you say that there are any shortcomings of the algorithm on this data set?
\end{itemize}

\textit{Note}: The runtime can be quite sensitive to the alphas of JCBB.
This is due both to the depth of search and creation of new landmarks, creating sort of a tradeoff.
This is especially true when other poor parameters are chosen, so that too few or too many landmarks fit the prediction.
This is usually seen quite quickly by the associations, and how many new landmarks are created.
It is also sensitive to \pythoninline{sensor.range_local} and \pythoninline{sensor.bearing_local_deg}, which decide how many variables the joint covariance in Task \ref{q:task01} (e) is recovered over.
